% Lösungen Einheit 1 · Klammern auflösen
\zweigkopf{Einheit 1 von 5 · Klammern auflösen}{Hier lernst du, Klammern mit Plus, Minus oder einer Zahl davor aufzulösen · neu oder schon bekannt – je nach Buch · keine P10-Aufgabe · baut auf: Terme zusammenfassen, Malnehmen mit Vorzeichen}{l1}
\erg{8}{a) $5x$, $-2y$, $+7$ \quad b) $-3m$, $+4$, $-n$ \quad c) $x^2$, $-6x$, $-1$ \quad d) $-8$, $+2x$, $-x^2$}
\erg{9}{a) Plus \quad b) Minus \quad c) positive Zahl \quad d) negative Zahl ($-4$) \quad e) Minus}
\erg{10}{a) $x + 7$ \quad b) $3x + 5$ \quad c) $x + 4$ \quad d) $3m + 2$}
\erg{11}{a) $2x + 6$ \quad b) $5m + 20$ \quad c) $6y + 6$ \quad d) $14 + 7x$ \quad e) $4x - 20$ \quad f) $2x + 18$ \quad g) $x^2 + 4x$}
\erg{12}{a) $-x - 6$ \quad b) $8 - x$ \quad c) $12 - x$ \quad d) $3x + 1$ \quad e) $3m + 2n - 3$ \quad f) $-3x + 12$ \quad g) $2x - 6$ \quad h) $6x + 10 - 4x + 8 = 2x + 18$}
\erg{13}{a) Jonas hat nur das erste Vorzeichen in der Klammer gedreht; richtig: $12 - x + 5 = 17 - x$ \quad b) $13 - 2x$ \quad c) Nina hat $5 - 2$ zuerst gerechnet (Punkt vor Strich missachtet); richtig: $5 - 2x - 2 = 3 - 2x$ \quad d) $2 - 3x$}
\erg{14}{a) nein: $3 \cdot (x - 2) = 3x - 6$, nicht $3x - 2$ \quad b) ja: $6 - x - 2 = 4 - x$ \quad c) nein: ein Wert reicht nicht; für $x = 0$ ergibt sich $2$ und $3$ ($2x + 2$ ist nicht $x + 3$)}
\erg{15}{a) $20 - 3 \cdot (p + 2)$ \quad b) $20 - 3p - 6 = 14 - 3p$ \quad c) $14 - 3 \cdot 1{,}5 = 9{,}50$~€ \quad d) nein: $14 - 3 \cdot 5 = -1$, es fehlt $1$~€}
